\documentclass{article}
\usepackage{amsmath}
\usepackage{amssymb}

\title{Mathe 3 Übungsblatt 8}
\author{Pascal Diller, Timo Rieke}
\date{}

\begin{document}
\maketitle

\section*{Aufgabe 1}
Nach mehrdimensionaler Kettenregel:
\[\frac{d}{dt}(f \circ \gamma)(t) = \Delta f(\gamma (t)) \cdot \gamma ' (t)\]
\[\text{Gradient }\Delta f(x,y) = \begin{pmatrix} \partial_xf(x,y) \\ \partial_yf(x,y) \end{pmatrix} = \begin{pmatrix} -2x \\ -2y \end{pmatrix} \]
Einsetzen $t = 0$, $\gamma(0) = (1,1)$
\[\Delta f(1,1) = \begin{pmatrix} -2 \\ -2 \end{pmatrix}\]
Mit $\gamma '(0) = (-1, -1)$ ergibt sich:
\[\frac{d}{dt} (f \circ \gamma) (0) = \begin{pmatrix}
    -2 \\ -2
\end{pmatrix} \cdot \begin{pmatrix}
    -1 \\ -1
\end{pmatrix} = (-2)(-1) + (-2)(-1) = 2 + 2 = 4\]
Gesucht ist die Richtungsableitung im Punkt $x_0 = \gamma(0)$ in Richtung des Tangentenvektors $v = \gamma'(0)$. \\
Die Definition der Richtungsableitung $\partial_v f(x_0)$ ist das Skalarprodukt aus dem Gradienten an der Stelle $x_0$ und dem Richtungsvektor $v$. Das ist das Produkt $\Delta f (\gamma(0)) \cdot \gamma'(0)$ mit dem Ergebnis 4, wie eben gezeigt.

\section*{Aufgabe 2}
\subsection*{(i)}
\[\partial_x f(x,y) = 1 \cdot \sin (y - z) + x \cdot \cos(y-x) \cdot (-1) \] 
\[= \sin(y-x) - x \cos (y-x)\]
\[\partial_y f(x,y) = x \cdot \cos(y-x) \cdot 1 \] 
\[= x \cos (y-x)\]
\[\partial_x f(1,1) = 0 -1 \cdot 1 = -1\]
\[\partial_y f(1,1) = 1 \cdot 1 = 1\]
\[\Delta f (1,1) = \begin{pmatrix}
    -1 \\ 1
\end{pmatrix}\]
\[\partial_v f(1,1) = \begin{pmatrix}
    -1 \\ 1
\end{pmatrix} \cdot \begin{pmatrix}
    \cos t \\ \sin t
\end{pmatrix} = - \cos t + \sin t = \sin t - \cos t \]
\subsection*{(ii)}
Die Richtungsableitung ist maximal, wenn $v$ in Richtung des Gradienten $\Delta f(1,1) = (-1, 1)$ zeigt. Der Vektor $(-1, 1)$ liegt im 2. Quadranten und halbiert diesen, was einem Winkel von $135^{\circ} = \frac{3 \pi}{4}$ entspricht. Der maximale Wert ist die Laenge des Gradienten $||\Delta f|| = \sqrt{(-1)^2 + 1^2} = \sqrt{2}$ \\
\newline
$\implies$ Der Wert ist maximal fuer $t = \frac{3\pi}{4}$ mit dem Wert $\sqrt{2}$

\subsection*{(iii)}
\[D^2 f = \begin{pmatrix}
    f_{xx} & f_{xy} \\ f_{yx} & f_{yy}
\end{pmatrix}\]
\[f_x(x,y) = \sin(y-x) - x\cos(y-x)\]
\[f_y(x,y) = x\cos(y-x)\]
\newline
\[f_{xx} = \frac{\partial}{\partial x}(\sin(y-x) - x\cos(y-x))\]
\[= \cos(y-x)\cdot(-1) - \left[ 1 \cdot \cos(y-x) + x \cdot (-\sin(y-x)) \cdot (-1) \right]\]
\[= -\cos(y-x) - \cos(y-x) - x\sin(y-x)\]
\[= -2\cos(y-x) - x\sin(y-x)\]
\newline
\[f_{yy} = \frac{\partial}{\partial y}(x\cos(y-x)) = x \cdot (-\sin(y-x)) \cdot 1 = -x\sin(y-x)\]
\newline
\[f_{xy} = \frac{\partial}{\partial x}(x\cos(y-x))\]
\[= 1 \cdot \cos(y-x) + x \cdot (-\sin(y-x)) \cdot (-1)\]
\[= \cos(y-x) + x\sin(y-x)\]
\newline
\[D^2f(x,y) = \begin{pmatrix}
-2\cos(y-x) - x\sin(y-x) & \cos(y-x) + x\sin(y-x) \\
\cos(y-x) + x\sin(y-x) & -x\sin(y-x)
\end{pmatrix}\]

\end{document}
